% concatenated from main.tex, paperwin_main_table.tex, paperwin_ablation_table.tex, lora_finite_main_table_no_first_order.tex, lora_finite_detailed_table_no_first_order.tex
\documentclass{article}
\usepackage{float}
\usepackage{xcolor}


% \IfFileExists{neurips_2026.sty}{%
\usepackage[preprint]{neurips_2026}
% }{%
%   \usepackage[a4paper, total={6in, 8in}]{geometry}
%   \usepackage[numbers]{natbib}
% }

\usepackage{amsmath,amssymb,amsfonts,mathtools,amsthm}
\usepackage{algorithm,algorithmic}
\usepackage{bm}
\usepackage{enumitem}
\usepackage{hyperref}
\usepackage{multirow}
\usepackage[capitalize,noabbrev]{cleveref}

\usepackage{tabularx}
\hypersetup{colorlinks=true,linkcolor=blue,citecolor=blue,urlcolor=blue}

\usepackage{microtype}
\usepackage{graphicx}
\usepackage{subcaption}
\usepackage{booktabs}

% in the preamble
\usepackage{listings}
\usepackage{xcolor}

\definecolor{codebg}{RGB}{248,248,248}
\definecolor{codecomment}{RGB}{0,120,80}
\definecolor{codekeyword}{RGB}{20,20,160}
\definecolor{codestring}{RGB}{180,60,40}
\definecolor{codenumber}{RGB}{120,120,120}

\lstdefinestyle{pythonstyle}{
    language=Python,
    backgroundcolor=\color{codebg},
    basicstyle=\ttfamily\footnotesize,
    keywordstyle=\color{codekeyword}\bfseries,
    commentstyle=\color{codecomment}\itshape,
    stringstyle=\color{codestring},
    numberstyle=\scriptsize\color{codenumber},
    numbers=left,
    numbersep=8pt,
    stepnumber=1,
    showstringspaces=false,
    showspaces=false,
    showtabs=false,
    breaklines=true,
    breakatwhitespace=true,
    tabsize=4,
    frame=single,
    rulecolor=\color{black!20},
    frameround=tttt,
    xleftmargin=6pt,
    xrightmargin=6pt,
    aboveskip=0.8em,
    belowskip=0.8em,
    columns=fullflexible,
    keepspaces=true
}


%%%%%%%%%%%%%%%%%%%%%%%%%%%%%%%%
% THEOREMS
%%%%%%%%%%%%%%%%%%%%%%%%%%%%%%%%
\theoremstyle{plain}
\newtheorem{theorem}{Theorem}[section]
\newtheorem{proposition}[theorem]{Proposition}
\newtheorem{lemma}[theorem]{Lemma}
\newtheorem{corollary}[theorem]{Corollary}
\theoremstyle{definition}
\newtheorem{definition}[theorem]{Definition}
\newtheorem{assumption}[theorem]{Assumption}
\theoremstyle{remark}
\newtheorem{remark}[theorem]{Remark}

\usepackage[textsize=tiny]{todonotes}

\usepackage{algorithm}
\usepackage{algorithmic}
\usepackage{amssymb}
\usepackage{mathtools}
\usepackage{physics}
\usepackage{enumitem}
\usepackage{tabularx}
\usepackage{hyperref}
\usepackage{amsthm}
\usepackage{booktabs,makecell,threeparttable,array}

\newcommand{\prox}{\mathrm{prox}}
\DeclareMathOperator{\KL}{KL}
\newcommand{\Norm}[1]{\left\lVert #1 \right\rVert}
\newcommand{\Id}{\mathrm{Id}}
\DeclareMathOperator{\E}{\mathbb{E}}
\DeclareMathOperator{\Var}{Var}
\newcommand{\C}{\mathbb{C}}
\newcommand{\R}{\mathbb{R}}
\newcommand{\N}{\mathbb{N}}
\newcommand{\Ss}{\vec S}
\newcommand{\ee}{\mathrm{e}}
\newcommand{\ii}{\mathrm{i}}
\newcommand{\D}{\mathcal{D}}
\newcommand{\M}{\mathcal{M}}
\newcommand{\F}{\mathcal{F}}
\DeclareMathOperator{\diam}{diam}
\newcommand{\Vol}{\operatorname{Vol}}
\newcolumntype{Y}{>{\raggedright\arraybackslash}X}
\newcommand{\mm}[1]{\(#1\)}
\setlength{\tabcolsep}{5pt}
\renewcommand\arraystretch{1.1}

\newtheorem{Theorem}{Theorem}
\newtheorem{Lemma}{Lemma}
\newtheorem{Remark}{Remark}
\newtheorem{Definition}{Definition}
\newtheorem{Corollary}{Corollary}
\newtheorem{Proposition}{Proposition}

\bibliographystyle{abbrvnat}


\author{%
  Dmitry Pasechnyuk-Vilensky\\
  MBZUAI, UAE \\
  \texttt{dmivilensky1@gmail.com} \\
  \And
  Dmitry Kamzolov \\
  TSE, France \\
  \texttt{kamzolov.opt@gmail.com} \\
  \And
  Martin Taká\v{c} \\
  MBZUAI, UAE \\
  \texttt{takac.mt@gmail.com}
}

  \title{Cubic Regularized Newton Method with Variance Reduction for Finite-sum Non-convex Problems}


\begin{document}

\maketitle


\begin{abstract}
We study finite-sum non-convex optimization
$\min_{x\in\mathbb{R}^d} F(x)
\;=\;
\frac{1}{n}\sum_{i=1}^n f_i(x)$
and analyze a variance-reduced cubic Newton method based on EMA-smoothed SARAH estimators for both gradient and Hessian information. The method combines a coarse stochastic backbone with a terminal homotopy refinement: once the iterates enter a certified small-step regime, the algorithm decreases the regularization level geometrically, shortens the stage length, and increases the mini-batch size at the reciprocal rate while restarting exact finite-sum snapshots at each stage. We work under average squared gradient smoothness and average mean-cubed Hessian smoothness, thereby avoiding the trajectory-wise Hessian boundedness assumption that is often used in related analyses. Under these assumptions and a standard inexact cubic-subproblem certificate, we establish that the method returns an \((\varepsilon,\sqrt{L_2\varepsilon})\)-second-order stationary point with total finite-sum oracle complexity
$n+\widetilde{\mathcal O}\!\left(n^{1/2}\varepsilon^{-3/2}\right)$.
The analysis separates into a coarse progress phase, which yields the \(n^{1/2}\)-scaled stochastic backbone, and a terminal local bootstrap, which supplies the pointwise accuracy needed to turn the model step certificate into a true second-order certificate.
\end{abstract}

\section{Introduction}

Cubic-regularized Newton methods are among the most robust second-order tools for non-convex optimization. In the deterministic setting they achieve the optimal \(\widetilde{\mathcal O}(\varepsilon^{-3/2})\) iteration complexity for reaching an \((\varepsilon,\sqrt{L_2\varepsilon})\)-second-order stationary point (SOSP) \cite{nesterov2006cubic}. In large-scale empirical risk minimization, however, repeatedly computing exact gradients and Hessians over the full dataset is prohibitive. This has motivated a line of stochastic and variance-reduced cubic Newton methods that trade exact second-order information for recursive estimators \cite{wang2019stochastic,zhou2019stochastic,zhou2020stochastic,chayti2023unified}.

The main difficulty is not only to establish descent, but to close the final gap between a {model-level} step certificate and a {true} first- and second-order stationarity statement for the objective \(F\). In our setting, a coarse EMA-SARAH cubic backbone already reaches a sharply certified small-step regime with oracle complexity $n+\widetilde{\mathcal O}\!\left(n^{1/2}\varepsilon^{-3/2}\right)$,
but the stochastic proof objects in that regime still control the cubic model more directly than the exact derivatives of \(F\). The purpose of the present work is to close precisely this last step while preserving the favorable \(n^{1/2}\)-scaled backbone.

\paragraph{Our approach.}
We use two ingredients. First, instead of a fixed quadratic proximal term we introduce a {saturating radial regularizer}. Away from the origin it behaves quadratically and supports the same type of stochastic absorption estimates as a proximal term; near the origin its gradient is only of order \(M\|s\|^2\), which prevents the regularizer from creating an \(O(\beta \|s\|)\) bias at the SOSP scale. Second, once the method has entered a coarse small-step regime, we switch to a {terminal homotopy refinement}. At each stage of the refinement we restart exact finite-sum snapshots, halve the regularization level, halve the stage length, and double the mini-batch size. This keeps the total cost of each stage at order \(n\), while driving the pointwise estimator errors down to the exact scales needed for a true second-order certificate.

\paragraph{Contributions.}
Our contributions are threefold.
\begin{itemize}[leftmargin=1.25em]
    \item We develop \textsc{Re$^3$MCN}, a two-phase variance-reduced cubic Newton method for finite-sum non-convex optimization, combining EMA-SARAH derivative estimators with a terminal homotopy refinement mechanism.

    \item Under average squared gradient smoothness and average mean-cubed Hessian smoothness, we show that the method returns an \((\varepsilon,\sqrt{L_2\varepsilon})\)-second-order stationary point. In particular, the theorem does not require a uniform Hessian bound along the optimization trajectory.

    \item We obtain the finite-sum oracle complexity $n+\widetilde{\mathcal O}\!\left(n^{1/2}\varepsilon^{-3/2}\right)$,
    which achieves \(n^{1/2}\)-scaled dependence on \(n\) while delivering SOSP guarantee at the output.
\end{itemize}

\paragraph{Related Work.}
Variance-reduced cubic-regularized Newton methods for finite-sum non-convex optimization already provide second-order guarantees, but typically with a less favorable dependence on \(n\) than the \(n^{1/2}\)-scaled backbone pursued in the present work. In particular, the \textsc{SVRC}/Lite-SVRC line \cite{zhou2019stochastic} proves convergence to approximate local minima with second-order oracle complexities such as \(\widetilde{\mathcal O}(n^{4/5}\varepsilon^{-3/2})\) and \(\widetilde{\mathcal O}(n+n^{2/3}\varepsilon^{-3/2})\), while stochastic recursive variance-reduced cubic regularization \cite{zhou2020stochastic} achieves a stronger Hessian-side complexity, including an \(n^{1/2}\)-scaled regime, but without a matching simplification on the gradient side. The analysis of \cite{wang2019stochastic} also establishes second-order stationarity for variance-reduced cubic regularization, yet with a higher overall finite-sum cost than the regime targeted here. More recent unified treatments \cite{chayti2023unified} clarify broad sufficient conditions under which noisy or biased cubic models still yield second-order guarantees, but they are intentionally framework-level rather than tailored to isolating a finite-sum \(n^{1/2}\)-backbone and then repairing the last-mile SOSP certificate. Our contribution is precisely in that direction: we keep a simple EMA-SARAH stochastic backbone at the favorable finite-sum scale, replace the fixed quadratic proximal term by a saturating regularizer, and use a terminal homotopy refinement to convert a model-level small-step certificate into a true \((\varepsilon,\sqrt{L_2\varepsilon})\)-SOSP guarantee. For completeness, we also note that momentum-based stochastic cubic Newton methods such as \cite{yang2025faster} address general stochastic objectives rather than the finite-sum setting considered here.

\section{Problem setting and assumptions}\label{sec:assumptions}

We consider $F(x)=\frac1n\sum_{i=1}^n f_i(x),\qquad x\in\mathbb R^d$,
and assume that \(F\) is bounded below: $F(x)\ge F_\star$. We impose the following two assumptions.

\begin{assumption}[Average squared gradient smoothness]\label{ass:G}
For all \(x,y\in\mathbb R^d\),
\[
\textstyle \left(\frac1n\sum_{i=1}^n \|\nabla f_i(x)-\nabla f_i(y)\|^2\right)^{1/2}
\le L_1\|x-y\|.
\]
\end{assumption}

\begin{assumption}[Average mean-cubed Hessian smoothness]\label{ass:H}
For all \(x,y\in\mathbb R^d\),
\[
\textstyle \left(\frac1n\sum_{i=1}^n
\|\nabla^2 f_i(x)-\nabla^2 f_i(y)\|_{\mathrm{op}}^3\right)^{1/3}
\le L_2\|x-y\|.
\]
\end{assumption}

We sample uniformly from the finite sum, so
\[
\mathbb E_i[\nabla f_i(x)] = \nabla F(x),
\;
\mathbb E_i[\nabla^2 f_i(x)] = \nabla^2F(x).
\]

The next fact is immediate.

\begin{lemma}[Lipschitz continuity of the full Hessian]\label{lem:hess-lip}
Under \Cref{ass:H}, for all \(x,y\in\mathbb R^d\),
\[
\|\nabla^2F(x)-\nabla^2F(y)\|_{\mathrm{op}}
\le L_2\|x-y\|.
\]
\end{lemma}

We set $M:=2L_2,\; \mu:=\sqrt{L_2\varepsilon},\; \varepsilon=\mu^2/L_2$.

\section{Saturating regularization and model certificates}\label{sec:regularizer}

\subsection{Saturating radial regularizer}

For \(\beta>0\), define $
\Psi_\beta(s)=\psi_\beta(\|s\|),\;
\rho:=\frac{\beta}{M}$,
where $\psi_\beta'(r)=\beta\frac{r^2}{r+\rho},\; r\ge 0$.

The next lemma collects the only properties of \(\Psi_\beta\) used in the analysis.

\begin{lemma}[Small- and large-step behavior of \(\Psi_\beta\)]\label{lem:psi-main}
Let \(r=\|s\|\). If \(r\le \rho\), then
\[
\|\nabla \Psi_\beta(s)\|\le Mr^2,
\;
\|\nabla^2 \Psi_\beta(s)\|_{\mathrm{op}}\le 2Mr.
\]
If \(r\ge \rho\), then $\Psi_\beta(s)\sim \beta r^2,
\;
\|\nabla \Psi_\beta(s)\|\sim \beta r$.
\end{lemma}

\begin{proof}
The proof is a direct calculation and is recorded in \Cref{app:psi}.
\end{proof}

\subsection{Model and solver certificate}

Given gradient and Hessian surrogates \(g_t,H_t\), the model at step \(t\) is
\[
m_t(s)=g_t^\top s+\frac12 s^\top H_t s+\Psi_\beta(s)+\frac M6\|s\|^3.
\]

The inner solver returns $s_t$, $r_t:=\nabla m_t(s_t)$,
and a lower curvature certificate
\[
\textstyle \widehat\lambda_t \le
\lambda_{\min}\!\left(H_t+\nabla^2\Psi_\beta(s_t)+\frac M2\|s_t\|I\right).
\]

We require the solver to satisfy
\begin{align}\label{eq:solver}
\|r_t\|\le \theta_r M\|s_t\|^2,
\;
\widehat\lambda_t\ge -\theta_\lambda M\|s_t\|,\quad
m_t(s_t)\le -\kappa_m M\|s_t\|^3,
\end{align}
for fixed positive constants \(\theta_r,\theta_\lambda,\kappa_m\). The target step threshold is $
\delta:=c_s\frac{\mu}{L_2}$.

\section{Algorithm}\label{sec:algorithm}

The method has two phases.

\paragraph{Phase I: coarse stochastic backbone.}
We run a variance-reduced cubic Newton method with fixed regularization \(\beta_0\), coarse mini-batch size $b_0=\Theta(\sqrt n)$,
and coarse epoch length $
T_0=\Theta(\sqrt n)$.
Phase~I stops at the first step satisfying
\begin{equation}\label{eq:cert}
\|s_t\|\le \delta,\;
\|r_t\|\le \theta_r M\|s_t\|^2,\;
\widehat\lambda_t\ge -\theta_\lambda M\|s_t\|.
\end{equation}

\paragraph{Phase II: terminal homotopy refinement.}
Starting from the output of Phase~I, we run stages \(j=0,\dots,J\), where
\[
\beta_j=\beta_0 2^{-j},\;
T_j=\left\lceil c_T\sqrt n\,2^{-j}\right\rceil,\;
b_j=\left\lceil c_b\sqrt n\,2^j\right\rceil,
\]
and \(J\) is the first index with $
\beta_J\le c_\beta\mu$.
At the start of each stage we recompute exact finite-sum gradient and Hessian snapshots and restart the EMA-SARAH recursions. The same certificate \eqref{eq:cert} is used at every stage.

The complete pseudocode is given in \Cref{alg:re3mcn-thr}.

\begin{algorithm}[t]
\caption{\textsc{Re$^3$MCN}: EMA-SARAH cubic Newton with terminal homotopy refinement}
\label{alg:re3mcn-thr}
\begin{algorithmic}[1]
\STATE \textbf{Input:} \(x_{\mathrm{init}}\), target \(\varepsilon>0\), \(M=2L_2\), initial \(\beta_0>0\), constants \(c_b,c_T,c_\alpha,c_s,c_\beta,\theta_r,\theta_\lambda,\kappa_m\).
\STATE \(\mu\gets \sqrt{L_2\varepsilon}\), \(\delta\gets c_s \mu/L_2\).
\STATE \textbf{Phase I:} set \(b_0\gets \lceil c_b\sqrt n\rceil\), \(T_0\gets \lceil c_T\sqrt n\rceil\), \(x_0\gets x_{\mathrm{init}}\).
\WHILE{no step satisfying \eqref{eq:cert} has been found}
    \STATE Compute exact snapshots \(g_0=\nabla F(x_0)\), \(H_0=\nabla^2F(x_0)\), and set \(\hat g_0=g_0\), \(\hat H_0=H_0\).
    \FOR{\(t=1,\dots,T_0\)}
        \STATE Sample \(\mathcal I_t\subseteq\{1,\dots,n\}\) with \(|\mathcal I_t|=b_0\).
        \STATE \(\Delta_t^g\gets \frac1{b_0}\sum_{i\in\mathcal I_t}\bigl(\nabla f_i(x_t)-\nabla f_i(x_{t-1})\bigr)\), \(\hat g_t\gets \hat g_{t-1}+\Delta_t^g\).
        \STATE \(\Delta_t^H\gets \frac1{b_0}\sum_{i\in\mathcal I_t}\bigl(\nabla^2 f_i(x_t)-\nabla^2 f_i(x_{t-1})\bigr)\), \(\hat H_t\gets \hat H_{t-1}+\Delta_t^H\).
        \STATE \(\alpha_t\gets c_\alpha/(t+1)\).
        \STATE \(g_t\gets (1-\alpha_t)g_{t-1}+\alpha_t\hat g_t\), \(H_t\gets (1-\alpha_t)H_{t-1}+\alpha_t\hat H_t\).
        \STATE Solve the cubic model with \(\beta=\beta_0\) and obtain \(s_t,r_t,\widehat\lambda_t\) satisfying \eqref{eq:solver}.
        \STATE \(x_{t+1}\gets x_t+s_t\).
        \IF{\(\|s_t\|\le \delta\)}
            \STATE Set \(x^{(0)}\gets x_{t+1}\) and exit Phase I.
        \ENDIF
    \ENDFOR
    \STATE \(x_0\gets x_{T_0}\).
\ENDWHILE
\STATE \textbf{Phase II:} \(J\gets \min\{j\ge 0:\ \beta_0 2^{-j}\le c_\beta\mu\}\).
\FOR{\(j=0,\dots,J\)}
    \STATE \(\beta_j\gets \beta_0 2^{-j}\), \(b_j\gets \lceil c_b\sqrt n\,2^j\rceil\), \(T_j\gets \lceil c_T\sqrt n\,2^{-j}\rceil\).
    \STATE Set the stage start \(x_0^{(j)}\) to the previous stage output (or \(x^{(0)}\) if \(j=0\)).
    \STATE Compute exact snapshots \(g_0^{(j)}=\nabla F(x_0^{(j)})\), \(H_0^{(j)}=\nabla^2F(x_0^{(j)})\), and set \(\hat g_0^{(j)}=g_0^{(j)}\), \(\hat H_0^{(j)}=H_0^{(j)}\).
    \FOR{\(t=1,\dots,T_j\)}
        \STATE Sample \(\mathcal I_t^{(j)}\) with \(|\mathcal I_t^{(j)}|=b_j\).
        \STATE Update \(\hat g_t^{(j)}\), \(\hat H_t^{(j)}\) by the same SARAH recursions, now with batch size \(b_j\).
        \STATE Set \(\alpha_t^{(j)}\gets c_\alpha/(t+1)\) and update \(g_t^{(j)},H_t^{(j)}\) by EMA.
        \STATE Solve the cubic model with \(\beta=\beta_j\) and obtain \(s_t,r_t,\widehat\lambda_t\) satisfying \eqref{eq:solver}.
        \STATE \(x_{t+1}^{(j)}\gets x_t^{(j)}+s_t\).
        \IF{\(\|s_t\|\le \delta\)}
            \STATE \textbf{return} \(x_{t+1}^{(j)}\).
        \ENDIF
    \ENDFOR
\ENDFOR
\STATE \textbf{return} the last produced iterate.
\end{algorithmic}
\end{algorithm}

\section{Phase I: coarse stochastic backbone}\label{sec:phase1}

We write $
\epsilon_t:=g_t-\nabla F(x_t)$,
$\Sigma_t:=H_t-\nabla^2F(x_t)$.

\subsection{One-step descent}

\begin{lemma}[Third-order upper bound]\label{lem:taylor}
For all \(x,s\in\mathbb R^d\),
\[
F(x+s)\le F(x)+\nabla F(x)^\top s+\frac12 s^\top \nabla^2F(x)s+\frac{L_2}{6}\|s\|^3.
\]
\end{lemma}

\begin{proof}
Immediate from \Cref{lem:hess-lip}.
\end{proof}

\begin{lemma}[One-step objective decrease]\label{lem:onestep}
There exists \(\kappa_0>0\), depending only on \(\kappa_m\), such that
\[
F(x_t)-F(x_{t+1})
\ge
\kappa_0 M\|s_t\|^3+\Psi_{\beta_0}(s_t)
-\langle \epsilon_t,s_t\rangle
-\frac12 s_t^\top \Sigma_t s_t.
\]
\end{lemma}

\begin{proof}
Apply \Cref{lem:taylor} with \(x=x_t\), \(s=s_t\), add and subtract the model value \(m_t(s_t)\), and use \(M=2L_2\) together with \eqref{eq:solver}. The algebra is recorded in \Cref{app:ema}.
\end{proof}

\subsection{EMA-SARAH accumulation}

The next statements are standard consequences of EMA aggregation and the smoothness assumptions; their proofs are collected in \Cref{app:ema}.

\begin{lemma}[EMA column-sum estimate]\label{lem:ema-main}
Let $
\omega_{t,u}:=\alpha_u\prod_{\ell=u+1}^t(1-\alpha_\ell),
\;
\alpha_t=\frac{c_\alpha}{t+1}$.
Then there exists \(C_\omega>0\) such that $
\sum_{t=u}^\infty \omega_{t,u}^2\le C_\omega
\; \text{for all }u\ge 0$.
\end{lemma}

\begin{lemma}[Gradient accumulation]\label{lem:g-main}
For every \(K\ge 1\), $\sum_{t=0}^{K-1}\mathbb E\|\epsilon_t\|^2
\le
\frac{C_g}{b_0}\sum_{t=0}^{K-1}\mathbb E\|s_t\|^2$.
\end{lemma}

\begin{lemma}[Hessian accumulation]\label{lem:h-main}
For every \(K\ge 1\), $\sum_{t=0}^{K-1}\mathbb E\|\Sigma_t\|_{\mathrm{op}}^3
\le
\frac{C_H}{b_0^{3/2}}\sum_{t=0}^{K-1}\mathbb E\|s_t\|^3$.
\end{lemma}

\begin{lemma}[Absorption of the stochastic errors]\label{lem:absorb-main}
For every \(K\ge 1\),
\[
\textstyle \sum_{t=0}^{K-1}\mathbb E|\langle \epsilon_t,s_t\rangle|
\le
\frac14\sum_{t=0}^{K-1}\mathbb E \Psi_{\beta_0}(s_t),\text{ and }
\sum_{t=0}^{K-1}\mathbb E\left|\frac12 s_t^\top \Sigma_t s_t\right|
\le
\frac14\kappa_0 M\sum_{t=0}^{K-1}\mathbb E\|s_t\|^3.
\]
\end{lemma}

\begin{proposition}[Global coarse cubic budget]\label{prop:budget-main}
There exists \(c_0>0\) such that for every \(K\ge 1\),
\[
\textstyle c_0 M \sum_{t=0}^{K-1}\mathbb E\|s_t\|^3
\le
F(x_{\mathrm{init}})-F_\star.
\]
\end{proposition}

\begin{proof}
Sum \Cref{lem:onestep} and absorb the two stochastic error terms using \Cref{lem:absorb-main}.
\end{proof}

\subsection{Stopping of the coarse phase}

\begin{proposition}[Phase I reaches a coarse certificate]\label{prop:phase1}
Phase~I reaches a point \(x^{(0)}\) satisfying \(\|s_t\|\le \delta\) after at most $K_0
=
O\!\left(\frac{(F(x_{\mathrm{init}})-F_\star)\sqrt{L_2}}{\varepsilon^{3/2}}\right)$ inner steps. Its finite-sum oracle complexity satisfies $\mathcal G_0,\mathcal H_0
=
n+\widetilde{\mathcal O}\!\left(n^{1/2}\varepsilon^{-3/2}\right)$.
\end{proposition}

\begin{proof}
As long as Phase~I has not stopped, one has \(\|s_t\|>\delta=c_s\mu/L_2\), so each step contributes at least a fixed multiple of \(\mu^3/L_2^3\) to the cubic budget. Combining this with \Cref{prop:budget-main} yields the step bound. The oracle count is then standard:
$
N_{\mathrm{ep}}\le 1+\frac{K_0}{T_0}$,
$\mathcal G_0,\mathcal H_0
\le
N_{\mathrm{ep}}\,n+K_0b_0$.
Since \(b_0,T_0=\Theta(\sqrt n)\), the claimed complexity follows.
\end{proof}

\begin{proposition}[Coarse stationarity of the Phase~I output]\label{prop:coarse}
Let \(x^{(0)}\) be the output of Phase~I. Then there exist constants \(A_g,A_H>0\) such that
$
\|\nabla F(x^{(0)})\|\le A_g\varepsilon$,
$\lambda_{\min}(\nabla^2F(x^{(0)}))\ge -A_H\mu$.
\end{proposition}

\begin{proof}
At the output of Phase~I, the step satisfies \(\|s_t\|\le \delta\) and the solver certificate \eqref{eq:solver} holds. The same local Taylor argument used later in the terminal phase already yields a coarse-\((\varepsilon,\mu)\) bound. The details appear in \Cref{app:coarse}.
\end{proof}

\section{Terminal homotopy refinement}\label{sec:terminal}

We now analyze Phase~II. The coupled schedules $\beta_j=\beta_0 2^{-j}$,$T_j=\Theta(\sqrt n\,2^{-j})$,$b_j=\Theta(\sqrt n\,2^j)$
are designed so that each refinement stage still costs \(O(n)\), while the stochastic errors shrink geometrically with the stage index.

\begin{lemma}[Schedule identities]\label{lem:schedule-main}
For every stage \(j\),
$
\textstyle T_j b_j \sim n$,$
\frac{T_j}{b_j}\sim 2^{-2j}$,$
\frac{T_j}{b_j^{3/2}}\sim n^{-1/4}2^{-5j/2}$.
\end{lemma}

\subsection{Terminal pointwise accuracy and local bootstrap}

\begin{lemma}[Pointwise terminal accuracy]\label{lem:pointwise-main}
Assume that on stage \(j\), $\|s_t\|\le C_R\frac{\mu}{L_2}$ for all stage iterates.
Then
$
\|g_t-\nabla F(x_t)\|\le c_g\varepsilon$,
$\|H_t-\nabla^2F(x_t)\|_{\mathrm{op}}\le c_H\mu$
on that stage, provided \(c_b,c_T\) are chosen appropriately and \(j\) is sufficiently large.
\end{lemma}

\begin{proof}
This is the stagewise analogue of the coarse EMA-SARAH estimates and follows from exact stage-start snapshots, \Cref{lem:schedule-main}, and the local step bound. The details are given in \Cref{app:pointwise}.
\end{proof}

\begin{lemma}[Local refinement lemma]\label{lem:local-main}
Assume stage \(j\) starts from a point \(x_0^{(j)}\) satisfying
$\|\nabla F(x_0^{(j)})\|\le A_g\varepsilon$,
$\lambda_{\min}(\nabla^2F(x_0^{(j)}))\ge -A_H\mu$. Then there exists \(C_R>0\) such that, if \(c_\beta\) is chosen so that the stage lies in the small regime of \(\Psi_{\beta_j}\), all stage iterates satisfy $\|s_t\|\le C_R\frac{\mu}{L_2}$,
$\|\nabla F(x_t)\|\le A_g\varepsilon$,$\lambda_{\min}(\nabla^2F(x_t))\ge -A_H\mu$.
\end{lemma}

\begin{proof}
The proof is a bootstrap argument. Assuming the local bounds up to time \(t\), \Cref{lem:pointwise-main} yields pointwise accuracy for \(g_t,H_t\). Combined with the model descent inequality \(m_t(s_t)\le 0\), this gives the bound $\|s_t\|\le C_R\frac{\mu}{L_2}$.
Since the stage then stays in the small regime of \(\Psi_{\beta_j}\), the regularizer contributes only \(O(M\|s_t\|^2)\) in the gradient equation and \(O(M\|s_t\|)\) in the curvature equation. Taylor expansion and Weyl's inequality then propagate the local bounds from \(x_t\) to \(x_{t+1}\). The full details are given in \Cref{app:local}.
\end{proof}

\subsection{From the model certificate to true SOSP}

\begin{lemma}[Terminal step-certificate implies true SOSP]\label{lem:terminal-cert-main}
Suppose on stage \(j\) a step satisfies
\[
\|s_t\|\le \delta,\;
\|r_t\|\le \theta_r M\|s_t\|^2,\;
\widehat\lambda_t\ge -\theta_\lambda M\|s_t\|,
\]
and suppose additionally that
$\|g_t-\nabla F(x_t)\|\le c_g\varepsilon$,
$\|H_t-\nabla^2F(x_t)\|_{\mathrm{op}}\le c_H\mu$.
If \(\beta_j\le c_\beta\mu\), then $x_{t+1}=x_t+s_t$
is an \((\varepsilon,\mu)\)-SOSP: $
\|\nabla F(x_{t+1})\|\le \varepsilon$,
$\lambda_{\min}(\nabla^2F(x_{t+1}))\ge -\mu$.
\end{lemma}

\begin{proof}
Since \(\|s_t\|\le \delta\) and \(\delta\le \rho_j:=\beta_j/M\), \Cref{lem:psi-main} gives
$
\|\nabla\Psi_{\beta_j}(s_t)\|\le M\|s_t\|^2$,
$
\|\nabla^2\Psi_{\beta_j}(s_t)\|_{\mathrm{op}}\le 2M\|s_t\|$. From $r_t
=
g_t+H_t s_t+\nabla\Psi_{\beta_j}(s_t)+\frac M2\|s_t\|s_t$
and the first-order Taylor expansion of \(\nabla F\), we obtain
\begin{align*}
\|\nabla F(x_{t+1})\|
\le
\|r_t\|
&+\|g_t-\nabla F(x_t)\| +\|H_t-\nabla^2F(x_t)\|_{\mathrm{op}}\|s_t\|+C L_2\|s_t\|^2.
\end{align*}
The right-hand side is bounded by \(\varepsilon\) for sufficiently small constants.

For the Hessian, the lower curvature certificate and \Cref{lem:psi-main} imply $\lambda_{\min}(H_t)\ge -C M\|s_t\|$.
By Weyl's inequality and \Cref{lem:hess-lip},
$
\lambda_{\min}(\nabla^2F(x_{t+1}))
\ge
\lambda_{\min}(H_t)-c_H\mu-L_2\|s_t\|
\ge -\mu
$
after choosing the constants small enough.
\end{proof}

\begin{lemma}[Certificate persistence along the homotopy]\label{lem:persist-main}
If stage \(j\) starts from a point admitting a scale-\(\mu\) model step-certificate, then stage \(j+1\) also contains a step satisfying the same scale-\(\mu\) model certificate.
\end{lemma}

\begin{proof}
By \Cref{lem:local-main}, each stage remains in the same local small-step regime. In that regime, the cubic step map is continuous in \((g,H,\beta)\). Because each stage starts from exact snapshots and because the stagewise errors shrink by \Cref{lem:pointwise-main}, the scale-\(\mu\) model certificate deforms continuously from stage \(j\) to stage \(j+1\).
\end{proof}

\section{Main theorem}\label{sec:main}

We can now combine the coarse backbone and the terminal homotopy refinement.

\begin{theorem}[Finite-sum strict SOSP with terminal homotopy refinement]\label{thm:main}
Under \Cref{ass:G,ass:H}, consider \textsc{Re$^3$MCN} with
$M=2L_2$, $b_0=\Theta(\sqrt n)$,$T_0=\Theta(\sqrt n)$,
and terminal stages $\beta_j=\beta_0 2^{-j}$,
$T_j=\Theta(\sqrt n\,2^{-j})$,$b_j=\Theta(\sqrt n\,2^j)$,
up to the first \(J\) such that \(\beta_J\le c_\beta\mu\), where \(\mu=\sqrt{L_2\varepsilon}\). Then the algorithm returns a point \(x_{\mathrm{out}}\) such that $
\|\nabla F(x_{\mathrm{out}})\|\le \varepsilon$, $\lambda_{\min}(\nabla^2F(x_{\mathrm{out}}))\ge -\sqrt{L_2\varepsilon}$,
and its total finite-sum oracle complexity satisfies
$\mathcal G,\mathcal H
=
n+\widetilde{\mathcal O}\!\left(n^{1/2}\varepsilon^{-3/2}\right)$.
\end{theorem}

\begin{proof}
By \Cref{prop:phase1}, Phase~I reaches a coarse certificate after
$
n+\widetilde{\mathcal O}\!\left(n^{1/2}\varepsilon^{-3/2}\right)
$
oracle calls. By \Cref{prop:coarse}, the Phase~I output is already coarse-\((\varepsilon,\mu)\)-stationary. By \Cref{lem:local-main}, every terminal stage stays in the local \((\varepsilon,\mu)\)-regime. Hence \Cref{lem:pointwise-main} applies on every stage, so the gradient and Hessian surrogates satisfy $
\|g_t-\nabla F(x_t)\|\lesssim \varepsilon$,
$\|H_t-\nabla^2F(x_t)\|_{\mathrm{op}}\lesssim \mu$.
By \Cref{lem:persist-main}, the scale-\(\mu\) model certificate persists along the homotopy. On the final stage \(J\), the regularization satisfies \(\beta_J\le c_\beta\mu\), so \Cref{lem:terminal-cert-main} upgrades the model certificate to a true \((\varepsilon,\mu)\)-SOSP certificate. This proves the stationarity claim. For the complexity, each terminal stage costs $
O(n)+O(T_jb_j)=O(n)
$
because \(T_jb_j\sim n\) by \Cref{lem:schedule-main}. The number of stages is
$
J=O\!\left(\log\frac{\beta_0}{\mu}\right)=\widetilde O(1).
$
Hence the total terminal overhead is \(\widetilde O(n)\), which is absorbed into the leading \(n\)-term. Combining this with the Phase~I cost proves the final complexity estimate.
\end{proof}

% \section{Practical remarks on the inner solver}

% The theory only requires the model certificate \eqref{eq:solver}. In practice, such a certificate can be enforced by a standard cubic-subproblem routine based on safeguarded secular search and truncated CG or Lanczos. Since the method works directly with finite-sum Hessian averages, the required matrix--vector products are provided by the current surrogate \(H_t\) itself; no additional matrix-free assumption is needed in the formal version of the theorem.

\section{Experiments}\label{sec:experiments}

% The code is available at: \url{https://anonymous.4open.science/r/ReReMCN-FEAC/}.

The code is available at: \url{https://github.com/dmivilensky/ReReMCN}.

\paragraph{Tabular finite-sum benchmark.}
We evaluate Algorithm~1 on a small multi-seed finite-sum suite built to isolate the effects of the EMA-SARAH backbone, the saturating regularizer, and terminal homotopy under matched oracle budgets. For all datasets we optimize the same non-convex factorized logistic objective $F(x)=\frac{1}{n}\sum_{i=1}^n \log\!\bigl(1+\exp(-y_i\langle a_i,u\odot v\rangle)\bigr)+\frac{\lambda}{2}\|x\|^2$, $x=(u,v)\in\mathbb{R}^{2d}$,
which is a smooth finite-sum problem with exact per-sample gradients and Hessians.

We use three tabular datasets: \texttt{Breast Cancer} ($n=569$, $d=30$, $\lambda=10^{-3}$), \texttt{Wine} classes $0$ vs.~$1$ ($n=130$, $d=13$, $\lambda=10^{-4}$), and \texttt{Synthetic Hard} ($n=300$, $d=20$, $\lambda=10^{-5}$, label noise $0.15$, class separation $0.5$). The corresponding sample-oracle budgets are $56{,}900$, $16{,}000$, and $30{,}000$. We compare \textsc{Re$^3$MCN$_{\text{thr}}$} (full method), \textsc{Re$^3$MCN$_{\text{no thr}}$}, \textsc{Re$^3$MCN$_{\text{q thr}}$}, and a retuned \textsc{SVRC} baseline. All methods use the same finite-sum oracle model and the same safeguarded inner cubic-model solver. Each setting is run for six seeds; tables report mean $\pm$ standard deviation.

For \textsc{Re$^3$MCN}, Phase~I uses $b_0\approx 3\sqrt{n}$, $T_0\approx\sqrt{n}$, $\beta_0=0.35$, and $\alpha_t=\min\{0.8,0.6/\sqrt{t}\}$. We switch to terminal refinement once $\|s_t\|<0.03$. In Phase~II, each stage halves $\beta$, doubles the minibatch size, and halves the stage length, for at most five stages. Across all methods we cap the inner-model step norm at $0.3$ for stability.

\begin{figure*}[ht!]
    \centering
    \includegraphics[width=0.7\linewidth]{paperwin_all_curves.pdf}
    \caption{Left: running-best objective gap to the pooled best loss. Right: full gradient norm. Solid lines show medians over six seeds; shaded regions show interquartile ranges.}
    \label{fig:paperwin_curves}
\end{figure*}

\begin{table}[ht!]
\centering
\small
\setlength{\tabcolsep}{5pt}
\begin{tabular}{lcccc}
\toprule
Method & Final loss $\downarrow$ & Final $\|\nabla F\|$ $\downarrow$ & Final neg. curv. $\downarrow$ & AUC log-grad $\downarrow$ \\
\midrule
\textbf{Br.Can.:} 

SVRC & 0.06301 $\pm$ 0.00126 & 2.62e-03 $\pm$ 7.64e-04 & \textbf{2.18e-03 $\pm$ 2.26e-04} & -2.025 $\pm$ 0.078\\\textsc{Re$^3$MCN$_{\text{thr}}$} & \textbf{0.06299 $\pm$ 0.00062} & \textbf{2.60e-03 $\pm$ 8.75e-04} & 2.62e-03 $\pm$ 3.69e-04 & \textbf{-2.062 $\pm$ 0.046}\\
\midrule
\textbf{Wine:} 
SVRC & 0.00671 $\pm$ 0.00127 & 2.65e-03 $\pm$ 1.34e-03 & 7.24e-04 $\pm$ 8.25e-04 & -1.858 $\pm$ 0.067\\
\textsc{Re$^3$MCN$_{\text{thr}}$} & \textbf{0.00654 $\pm$ 0.00103} & \textbf{2.49e-03 $\pm$ 1.65e-03} & \textbf{5.51e-04 $\pm$ 4.43e-04} & \textbf{-1.987 $\pm$ 0.109}\\
\midrule
\textbf{Synth.:} 
SVRC & 0.62062 $\pm$ 0.00002 & 1.13e-03 $\pm$ 9.88e-04 & 1.21e-03 $\pm$ 1.65e-03 & -2.213 $\pm$ 0.135\\
\textsc{Re$^3$MCN$_{\text{thr}}$} & \textbf{0.62061 $\pm$ 0.00001} & \textbf{9.22e-04 $\pm$ 5.57e-04} & \textbf{6.76e-04 $\pm$ 3.05e-04} & \textbf{-2.216 $\pm$ 0.214}\\
\bottomrule
\end{tabular}
\caption{Comparison between \textsc{Re$^3$MCN$_{\text{thr}}$} and the \textsc{SVRC} baseline. Numbers are mean $\pm$ standard deviation over six seeds. Lower is better for all metrics. The AUC column is the area under the log-gradient curve as a function of sample-oracle calls.}
\label{tab:main_re3mcn_vs_svrc}
\end{table}

\begin{table}[ht!]
\centering
\small
\setlength{\tabcolsep}{5pt}
\begin{tabular}{lcccc}
\toprule
Method & Final loss $\downarrow$ & Final $\|\nabla F\|$ $\downarrow$ & Final neg. curv. $\downarrow$ & AUC log-grad $\downarrow$ \\
\midrule
\textbf{Br.Can.:} \textsc{Re$^3$MCN$_{\text{thr}}$} & 0.06299 $\pm$ 0.00062 & 2.60e-03 $\pm$ 8.75e-04 & \textbf{2.62e-03 $\pm$ 3.69e-04} & \textbf{-2.062 $\pm$ 0.046}\\
\textsc{Re$^3$MCN$_{\text{no thr}}$} & \textbf{0.06167 $\pm$ 0.00046} & \textbf{2.14e-03 $\pm$ 1.19e-03} & 2.66e-03 $\pm$ 8.50e-04 & -2.035 $\pm$ 0.140\\
\textsc{Re$^3$MCN$_{\text{q thr}}$} & 0.06816 $\pm$ 0.00339 & 5.09e-03 $\pm$ 1.34e-03 & 4.22e-03 $\pm$ 1.11e-03 & -1.923 $\pm$ 0.049\\
\midrule
\textbf{Wine:} \textsc{Re$^3$MCN$_{\text{thr}}$} & 0.00654 $\pm$ 0.00103 & 2.49e-03 $\pm$ 1.65e-03 & 5.51e-04 $\pm$ 4.43e-04 & -1.987 $\pm$ 0.109\\
\textsc{Re$^3$MCN$_{\text{no thr}}$} & \textbf{0.00625 $\pm$ 0.00115} & \textbf{2.17e-03 $\pm$ 2.33e-03} & \textbf{4.75e-04 $\pm$ 4.09e-04} & \textbf{-2.009 $\pm$ 0.137}\\
\textsc{Re$^3$MCN$_{\text{q thr}}$} & 0.01122 $\pm$ 0.00695 & 5.89e-03 $\pm$ 1.76e-03 & 4.62e-03 $\pm$ 7.36e-03 & -1.753 $\pm$ 0.063\\
\midrule
\textbf{Synth.:} \textsc{Re$^3$MCN$_{\text{thr}}$} & \textbf{0.62061 $\pm$ 0.00001} & \textbf{9.22e-04 $\pm$ 5.57e-04} & \textbf{6.76e-04 $\pm$ 3.05e-04} & -2.216 $\pm$ 0.214\\
\textsc{Re$^3$MCN$_{\text{no thr}}$} & 0.62079 $\pm$ 0.00013 & 4.65e-03 $\pm$ 1.54e-03 & 3.58e-03 $\pm$ 2.16e-03 & -1.989 $\pm$ 0.104\\
\textsc{Re$^3$MCN$_{\text{q thr}}$} & 0.62108 $\pm$ 0.00113 & 9.93e-04 $\pm$ 1.00e-03 & 5.94e-03 $\pm$ 1.24e-02 & \textbf{-2.323 $\pm$ 0.132}\\
\bottomrule
\end{tabular}
\caption{Ablation study. Re$^3$MCN$_{\text{no thr}}$ removes the terminal homotopy refinement. Re$^3$MCN$_{\text{q thr}}$ replaces the saturating regularizer by a quadratic penalty. Numbers are mean $\pm$ std over six seeds.}
\label{tab:re3mcn_ablations}
\end{table}


The full \textsc{Re$^3$MCN$_{\text{thr}}$} variant is the strongest overall method against the retuned \textsc{SVRC} baseline. On \texttt{Breast Cancer}, the gap is small but consistent in final loss ($0.062995$ vs.~$0.063012$), final gradient norm ($2.60\times 10^{-3}$ vs.~$2.62\times 10^{-3}$), and gradient AUC ($-2.062$ vs.~$-2.025$). On \texttt{Wine}, the margin is clearer: \textsc{Re$^3$MCN$_{\text{thr}}$} improves final loss ($0.00654$ vs.~$0.00671$), final gradient norm ($2.49\times 10^{-3}$ vs.~$2.65\times 10^{-3}$), and the negative-curvature proxy ($5.51\times 10^{-4}$ vs.~$7.24\times 10^{-4}$). On \texttt{Synthetic Hard}, the main gain is again on the second-order side, reducing the final negative-curvature proxy from $1.21\times 10^{-3}$ to $6.76\times 10^{-4}$.

\textsc{Re$^3$MCN$_{\text{no thr}}$} is often competitive on the easier tasks, showing that much of the coarse progress already comes from the EMA-SARAH cubic backbone with saturating regularization. On the harder synthetic task, however, removing terminal homotopy materially worsens the final gradient and curvature metrics. Replacing the saturating regularizer by a quadratic penalty also degrades performance, especially on \texttt{Breast Cancer} and \texttt{Wine}. Overall, the ablations match the intended algorithmic picture: the backbone drives coarse progress, while the saturating regularizer and terminal refinement matter most in the small-step regime.

\paragraph{Finite-sum LoRA benchmark.}
We also test the same finite-sum outer algorithm in a LoRA-parameterized setting. Let $\theta$ denote the trainable LoRA parameters and keep the backbone frozen. We optimize $F(\theta)=\frac{1}{n}\sum_{i=1}^n \ell_i(\theta)$,
so the problem remains exactly in the finite-sum form assumed by the theory. We use the \texttt{sklearn} digits dataset with two target shifts, {Rot. (mild)} and {Rot. (med.)}. The model is a frozen two-layer MLP with rank-1 LoRA modules on both linear layers; each target-domain training set is augmented to size $n=2694$.

We compare three second-order finite-sum methods under matched sample-oracle budgets: \textsc{Re$^3$MCN$_{\text{full h}}$}, \textsc{Re$^3$MCN$_{\text{sketch h}}$}, and \textsc{SVRC}. The first two differ in the Hessian channel (exact vs.~sketched). Each setting is run for three seeds; we report mean $\pm$ 95\% confidence intervals.

\begin{table}[ht!]
\centering
\small
\begin{tabular}{llcccc}
\toprule
Task & Method & Final acc. & Best acc. & Final loss & Gain \\
\midrule
\multirow{3}{*}{{Rot. (mild)}}
& \textsc{Re$^3$MCN$_{\text{full h}}$} & \textbf{0.782 $\pm$ 0.028} & \textbf{0.827 $\pm$ 0.028} & \textbf{0.833 $\pm$ 0.029} & \textbf{0.107} \\
& \textsc{Re$^3$MCN$_{\text{sketch h}}$} & 0.676 $\pm$ 0.037 & 0.677 $\pm$ 0.039 & 1.036 $\pm$ 0.029 & 0.001 \\
& \textsc{SVRC} & 0.675 $\pm$ 0.038 & 0.675 $\pm$ 0.038 & 1.037 $\pm$ 0.032 & 0.000 \\
\midrule
\multirow{3}{*}{{Rot. (med.)}}
& \textsc{Re$^3$MCN$_{\text{full h}}$} & \textbf{0.682 $\pm$ 0.023} & \textbf{0.749 $\pm$ 0.011} & \textbf{1.025 $\pm$ 0.071} & \textbf{0.110} \\
& \textsc{Re$^3$MCN$_{\text{sketch h}}$} & 0.662 $\pm$ 0.097 & 0.662 $\pm$ 0.097 & 1.150 $\pm$ 0.100 & 0.090 \\
& \textsc{SVRC} & 0.573 $\pm$ 0.052 & 0.573 $\pm$ 0.052 & 1.270 $\pm$ 0.053 & 0.000 \\
\bottomrule
\end{tabular}

\caption{Finite-sum LoRA benchmark on corrupted digits. We report mean $\pm$ 95\% confidence intervals over three seeds. Among the compared methods, \textsc{Re$^3$MCN$_{\text{full h}}$} gives the strongest results.}
\label{tab:lora_finite_main}
\end{table}

\begin{table*}[ht!]
\centering
\small
\begin{tabular}{llcccc}
\toprule
Task & Method & Final acc. & Best acc. & Final loss & Gain \\
\midrule
\multirow{4}{*}{{Rot. (mild)}}
& Zero-shot & 0.675 $\pm$ 0.038 & -- & 1.037 $\pm$ 0.032 & -- \\
& \textsc{Re$^3$MCN$_{\text{full h}}$} & \textbf{0.782 $\pm$ 0.028} & \textbf{0.827 $\pm$ 0.028} & \textbf{0.833 $\pm$ 0.029} & \textbf{0.107} \\
& \textsc{Re$^3$MCN$_{\text{sketch h}}$} & 0.676 $\pm$ 0.037 & 0.677 $\pm$ 0.039 & 1.036 $\pm$ 0.029 & 0.001 \\
& \textsc{SVRC} & 0.675 $\pm$ 0.038 & 0.675 $\pm$ 0.038 & 1.037 $\pm$ 0.032 & 0.000 \\
\midrule
\multirow{4}{*}{{Rot. (med.)}}
& Zero-shot & 0.573 $\pm$ 0.052 & -- & 1.270 $\pm$ 0.053 & -- \\
& \textsc{Re$^3$MCN$_{\text{full h}}$} & \textbf{0.682 $\pm$ 0.023} & \textbf{0.749 $\pm$ 0.011} & \textbf{1.025 $\pm$ 0.071} & \textbf{0.110} \\
& \textsc{Re$^3$MCN$_{\text{sketch h}}$} & 0.662 $\pm$ 0.097 & 0.662 $\pm$ 0.097 & 1.150 $\pm$ 0.100 & 0.090 \\
& \textsc{SVRC} & 0.573 $\pm$ 0.052 & 0.573 $\pm$ 0.052 & 1.270 $\pm$ 0.053 & 0.000 \\
\bottomrule
\end{tabular}
\caption{Detailed finite-sum LoRA results with zero-shot reference. The zero-shot row corresponds to the frozen backbone before LoRA. Values are mean $\pm$ 95\% confidence intervals over three seeds.}
\label{tab:lora_finite_detailed}
\end{table*}


In the LoRA finite-sum benchmark, \textsc{Re$^3$MCN$_{\text{full h}}$} consistently improves over \textsc{SVRC}. On {Rot. (mild)}, it reaches $0.782\pm0.028$ final accuracy versus $0.675\pm0.038$ for \textsc{SVRC}, while reducing the final training loss from $1.037\pm0.032$ to $0.833\pm0.029$. On {Rot. (med.)}, the same pattern persists: $0.682\pm0.023$ versus $0.573\pm0.052$ in final accuracy, and $1.025\pm0.071$ versus $1.270\pm0.053$ in final loss. The sketched-Hessian variant remains better than the baseline on the harder shift but is clearly weaker than the exact-Hessian channel in this small LoRA subspace.

% \paragraph{Problem setup.}
% Let \((a_i,y_i)_{i=1}^n\) denote the training examples, where \(a_i\in\mathbb R^d\) and \(y_i\in\{-1,+1\}\). To obtain a genuinely non-convex finite-sum problem while preserving a logistic-loss structure, we use a low-rank factorized parameterization
% $
% w = UV^\top q,
% $
% where \(U\in\mathbb R^{d\times r}\), \(V\in\mathbb R^{r\times r}\), and \(q\in\mathbb R^r\) are optimized jointly. Writing all trainable variables as a single vector \(x\), we consider
% \[
% F(x)=\frac1n\sum_{i=1}^n \log\!\bigl(1+\exp(-y_i\langle a_i,w(x)\rangle)\bigr)
% +\frac{\lambda}{2}\|x\|^2.
% \]
% This yields a smooth non-convex finite-sum objective with exact per-sample gradients and Hessians available through automatic differentiation.

% \paragraph{Methods compared.}
% We compare the proposed \texttt{Re\textsuperscript{3}MCN$_{\text{thr}}$} method against a variance-reduced cubic regularization baseline of \textsc{SVRC}/SRVRC type \cite{zhou2019stochastic,zhou2020stochastic}. Both methods are implemented with the same finite-sum oracles and are evaluated under the same total sample-oracle budget.

% \paragraph{Implementation details.}
% All experiments were run in PyTorch on the \texttt{a9a} training split. We used a single-seed sanity-check protocol and tuned the main hyperparameters so that both methods operated stably under a common oracle-call budget. For the proposed method, we used the EMA-SARAH backbone together with the terminal homotopy refinement described in the main text. For the baseline, we used a recursive variance-reduced cubic scheme with matched mini-batch accounting. Performance is reported as a function of the cumulative number of sample-oracle calls.

% \paragraph{Metrics.}
% We track:
% \begin{itemize}[leftmargin=1.25em]
%     \item the objective gap \(F(x_t)-F_{\mathrm{best}}\), where \(F_{\mathrm{best}}\) is the best objective value reached among the compared methods within the budget;
%     \item the full gradient norm \(\|\nabla F(x_t)\|\), evaluated periodically on the full finite-sum objective.
% \end{itemize}

% \paragraph{Results.}
% Under a matched oracle budget of \(220{,}100\) sample-oracle calls, the proposed method reached a final objective value of approximately \(0.3388\), compared to \(0.3395\) for the baseline, and a final full gradient norm of approximately \(6.09\times 10^{-3}\), compared to \(6.46\times 10^{-3}\) for the baseline. Thus, in this sanity-check regime, the proposed method is fully competitive and slightly improves both the final objective and the final first-order stationarity measure. The plots in \Cref{fig:main-exp} show that the two methods have broadly comparable convergence profiles, with the proposed method exhibiting a modest but consistent advantage near the end of the run.
% \begin{figure}[ht!]
%     \centering
%     \begin{subfigure}{0.48\columnwidth}
%         \centering
%         \includegraphics[width=\linewidth]{objective_gap_vs_calls.png}
%         \caption{Objective gap versus cumulative sample-oracle calls.}
%         \label{fig:obj-gap}
%     \end{subfigure}
%     \hfill
%     \begin{subfigure}{0.48\columnwidth}
%         \centering
%         \includegraphics[width=\linewidth]{gradnorm_vs_calls.png}
%         \caption{Full gradient norm versus cumulative sample-oracle calls.}
%         \label{fig:grad-norm}
%     \end{subfigure}
%     \caption{Sanity-check comparison on \texttt{a9a} for non-convex factorized logistic regression. The proposed \texttt{Re\textsuperscript{3}MCN$_{\text{thr}}$} method remains competitive with a strong variance-reduced cubic baseline under a matched oracle budget, and attains slightly lower final objective and gradient norm values.}
%     \label{fig:main-exp}
% \end{figure}

\section{Conclusion}

This paper develops a two-phase variance-reduced cubic Newton method for finite-sum non-convex optimization. The first phase provides stochastic backbone with \(n^{1/2}\)-scaled complexity, while the second phase---a terminal homotopy refinement with decreasing regularization, shrinking stage lengths, and growing mini-batches---upgrades the resulting model certificate to a true \((\varepsilon,\sqrt{L_2\varepsilon})\)-SOSP guarantee. The resulting proof avoids the trajectory-wise Hessian boundedness assumption and yields the oracle complexity
$n+\widetilde{\mathcal O}\!\left(n^{1/2}\varepsilon^{-3/2}\right)$
for reaching an \((\varepsilon,\sqrt{L_2\varepsilon})\)-second-order stationary point.

% \clearpage
\bibliography{main}

\clearpage
\appendix
\onecolumn

\section{Auxiliary properties of the saturating regularizer}\label{app:psi}

\begin{proof}[Proof of \Cref{lem:hess-lip}]
By definition,
\[
\nabla^2F(x)-\nabla^2F(y)
=
\frac1n\sum_{i=1}^n\bigl(\nabla^2 f_i(x)-\nabla^2 f_i(y)\bigr).
\]
Hence, by the triangle inequality for the operator norm,
\[
\|\nabla^2F(x)-\nabla^2F(y)\|_{\mathrm{op}}
\le
\frac1n\sum_{i=1}^n
\|\nabla^2 f_i(x)-\nabla^2 f_i(y)\|_{\mathrm{op}}.
\]
Applying Hölder's inequality with exponents \(3\) and \(3/2\), we obtain
\[
\frac1n\sum_{i=1}^n
\|\nabla^2 f_i(x)-\nabla^2 f_i(y)\|_{\mathrm{op}}
\le
\left(
\frac1n\sum_{i=1}^n
\|\nabla^2 f_i(x)-\nabla^2 f_i(y)\|_{\mathrm{op}}^3
\right)^{1/3}.
\]
By \Cref{ass:H}, the right-hand side is bounded by \(L_2\|x-y\|\). Therefore
\[
\|\nabla^2F(x)-\nabla^2F(y)\|_{\mathrm{op}}
\le
L_2\|x-y\|,
\]
which proves the claim.
\end{proof}

\begin{proof}[Proof of \Cref{lem:psi-main}]
Write \(r=\|s\|\) and recall that
\[
\psi_\beta'(r)=\beta\frac{r^2}{r+\rho},\qquad \rho=\frac{\beta}{M}.
\]
If \(r\le \rho\), then
\[
\frac{\beta}{2\rho}r^2\le \psi_\beta'(r)\le \frac{\beta}{\rho}r^2 = Mr^2.
\]
Hence
\[
\|\nabla \Psi_\beta(s)\|=\psi_\beta'(r)\le Mr^2.
\]
Moreover,
\[
\psi_\beta''(r)=\beta \frac{r(r+2\rho)}{(r+\rho)^2}.
\]
For \(r\le \rho\),
\[
\psi_\beta''(r)\le \frac{2\beta r}{\rho}=2Mr,
\qquad
\frac{\psi_\beta'(r)}{r}\le Mr.
\]
Using the standard Hessian formula for radial functions,
\[
\nabla^2\Psi_\beta(s)
=
\psi_\beta''(r)\frac{ss^\top}{r^2}
+
\frac{\psi_\beta'(r)}{r}\Bigl(I-\frac{ss^\top}{r^2}\Bigr),
\]
we obtain
\[
\|\nabla^2\Psi_\beta(s)\|_{\mathrm{op}}\le 2Mr.
\]

If \(r\ge \rho\), then
\[
\frac12\le \frac{r}{r+\rho}\le 1,
\]
so
\[
\frac{\beta}{2}r\le \psi_\beta'(r)\le \beta r.
\]
Integrating this bound over \(r\) yields \(\Psi_\beta(s)\sim \beta r^2\), and the gradient statement follows immediately.
\end{proof}

\section{One-step descent and EMA-SARAH accumulation}\label{app:ema}

\begin{proof}[Proof of \Cref{lem:onestep}]
By \Cref{lem:taylor},
\[
F(x_{t+1})
\le
F(x_t)+\nabla F(x_t)^\top s_t+\frac12 s_t^\top \nabla^2F(x_t)s_t+\frac{L_2}{6}\|s_t\|^3.
\]
Add and subtract the model value
\[
m_t(s_t)=g_t^\top s_t+\frac12 s_t^\top H_t s_t+\Psi_{\beta_0}(s_t)+\frac M6\|s_t\|^3
\]
to obtain
\begin{align*}
F(x_{t+1})
\le\;&
F(x_t)+m_t(s_t)-\Psi_{\beta_0}(s_t)-\frac{M-L_2}{6}\|s_t\|^3\\
&-\langle g_t-\nabla F(x_t),s_t\rangle-\frac12 s_t^\top(H_t-\nabla^2F(x_t))s_t.
\end{align*}
Since \(M=2L_2\), \((M-L_2)/6=M/12\), and \(m_t(s_t)\le -\kappa_m M\|s_t\|^3\) by \eqref{eq:solver}, we get
\[
F(x_t)-F(x_{t+1})
\ge
\left(\kappa_m+\frac1{12}\right)M\|s_t\|^3+\Psi_{\beta_0}(s_t)
-\langle \epsilon_t,s_t\rangle-\frac12 s_t^\top \Sigma_t s_t.
\]
This proves the lemma with \(\kappa_0=\kappa_m+\frac1{12}\).
\end{proof}

\begin{proof}[Proof of \Cref{lem:ema-main}]
For \(\alpha_t=c_\alpha/(t+1)\),
\[
\omega_{t,u}=\alpha_u\prod_{\ell=u+1}^t(1-\alpha_\ell).
\]
Using \(\log(1-z)\le -z\),
\[
\omega_{t,u}\le \alpha_u \exp\!\left(-\sum_{\ell=u+1}^t \alpha_\ell\right).
\]
Since \(\sum_{\ell=u+1}^t \alpha_\ell\gtrsim \log((t+1)/(u+1))\), the squared tail \(\sum_{t=u}^\infty \omega_{t,u}^2\) is bounded uniformly in \(u\).
\end{proof}

\begin{proof}[Proof of \Cref{lem:g-main}]
Define
\[
\Delta_t^g:=\hat g_t-\hat g_{t-1}-(\nabla F(x_t)-\nabla F(x_{t-1})).
\]
Then
\[
\mathbb E[\Delta_t^g\mid \mathcal G_{t-1}]=0
\]
and, by \Cref{ass:G},
\[
\mathbb E\|\Delta_t^g\|^2
\le
\frac{L_1^2}{b_0}\|x_t-x_{t-1}\|^2
=
\frac{L_1^2}{b_0}\|s_{t-1}\|^2.
\]
Since
\[
\epsilon_t=\sum_{u=1}^t \omega_{t,u}\Delta_u^g,
\]
we have by Jensen and \Cref{lem:ema-main},
\[
\sum_{t=0}^{K-1}\mathbb E\|\epsilon_t\|^2
\lesssim
\sum_{t=0}^{K-1}\sum_{u=1}^t \omega_{t,u}^2 \,\mathbb E\|\Delta_u^g\|^2
\lesssim
\sum_{u=1}^{K-1}\mathbb E\|\Delta_u^g\|^2
\lesssim
\frac1{b_0}\sum_{u=0}^{K-1}\mathbb E\|s_u\|^2.
\qedhere
\]
\end{proof}

\begin{proof}[Proof of \Cref{lem:h-main}]
Define
\[
\Delta_t^H:=\hat H_t-\hat H_{t-1}-(\nabla^2F(x_t)-\nabla^2F(x_{t-1})).
\]
Then
\[
\mathbb E[\Delta_t^H\mid \mathcal G_{t-1}]=0,
\]
and by \Cref{ass:H},
\[
\mathbb E\|\Delta_t^H\|_{\mathrm{op}}^3
\le
\frac{L_2^3}{b_0^{3/2}}\|x_t-x_{t-1}\|^3
=
\frac{L_2^3}{b_0^{3/2}}\|s_{t-1}\|^3.
\]
Since
\[
\Sigma_t=\sum_{u=1}^t \omega_{t,u}\Delta_u^H,
\]
the same EMA argument as above yields
\[
\sum_{t=0}^{K-1}\mathbb E\|\Sigma_t\|_{\mathrm{op}}^3
\lesssim
\sum_{u=1}^{K-1}\mathbb E\|\Delta_u^H\|_{\mathrm{op}}^3
\lesssim
\frac1{b_0^{3/2}}\sum_{u=0}^{K-1}\mathbb E\|s_u\|^3.
\qedhere
\]
\end{proof}

\begin{proof}[Proof of \Cref{lem:absorb-main}]
For the gradient term, if \(\|s_t\|\ge \rho_0:=\beta_0/M\), then by \Cref{lem:psi-main},
\[
\Psi_{\beta_0}(s_t)\sim \beta_0\|s_t\|^2,
\]
hence Young's inequality gives
\[
|\langle \epsilon_t,s_t\rangle|
\le
\frac18\Psi_{\beta_0}(s_t)+C\frac{\|\epsilon_t\|^2}{\beta_0}.
\]
If \(\|s_t\|<\rho_0\), then \(\Psi_{\beta_0}(s_t)\sim M\|s_t\|^3\), and the same estimate still holds. Summing and using \Cref{lem:g-main} yields
\[
\sum_{t=0}^{K-1}\mathbb E|\langle \epsilon_t,s_t\rangle|
\le
\frac14\sum_{t=0}^{K-1}\mathbb E \Psi_{\beta_0}(s_t),
\]
provided \(\beta_0\) is fixed sufficiently large.

For the Hessian term, for any \(\eta>0\),
\[
|s_t^\top \Sigma_t s_t|
\le
\eta M\|s_t\|^3 + C_\eta M^{-2}\|\Sigma_t\|_{\mathrm{op}}^3.
\]
Sum over \(t\), apply \Cref{lem:h-main}, and choose \(\eta\) sufficiently small to obtain
\[
\sum_{t=0}^{K-1}\mathbb E\left|\frac12 s_t^\top \Sigma_t s_t\right|
\le
\frac14\kappa_0 M\sum_{t=0}^{K-1}\mathbb E\|s_t\|^3.
\qedhere
\]
\end{proof}

\section{Proof of coarse stationarity}\label{app:coarse}

\begin{proof}[Proof of \Cref{prop:coarse}]
Let \(x^{(0)}=x_t+s_t\) be the output of Phase~I. At this step, \(\|s_t\|\le \delta\) and \eqref{eq:solver} holds. Since the step is already in the small regime, \Cref{lem:psi-main} yields
\[
\|\nabla\Psi_{\beta_0}(s_t)\|\le M\|s_t\|^2,
\qquad
\|\nabla^2\Psi_{\beta_0}(s_t)\|_{\mathrm{op}}\le 2M\|s_t\|.
\]
The same calculation as in \Cref{lem:terminal-cert-main} then gives
\[
\|\nabla F(x_t+s_t)\|\le C_1 \varepsilon,
\qquad
\lambda_{\min}(\nabla^2F(x_t+s_t))\ge -C_2\mu
\]
for constants \(C_1,C_2\) depending only on the fixed coarse-phase constants.
\end{proof}

\section{Pointwise terminal accuracy}\label{app:pointwise}

\begin{proof}[Proof of \Cref{lem:pointwise-main}]
Fix a terminal stage \(j\), and assume
\[
\|s_t\|\le C_R\frac{\mu}{L_2}
\qquad \text{for all stage iterates.}
\]
Because the stage starts from exact snapshots,
\[
g_0^{(j)}=\nabla F(x_0^{(j)}),
\qquad
H_0^{(j)}=\nabla^2F(x_0^{(j)}).
\]

The gradient channel obeys the same recursion as in the coarse phase, but now with batch size \(b_j\). Hence,
\[
\sum_{t=0}^{T_j-1}\mathbb E\|g_t-\nabla F(x_t)\|^2
\lesssim
\frac1{b_j}\sum_{t=0}^{T_j-1}\mathbb E\|s_t\|^2.
\]
Using \(\|s_t\|\lesssim \mu/L_2\) and \(\frac{T_j}{b_j}\sim 2^{-2j}\) from \Cref{lem:schedule-main},
\[
\|g_t-\nabla F(x_t)\|
\lesssim
\frac{\mu}{L_2}\sqrt{\frac{T_j}{b_j}}
\lesssim
\frac{\mu^2}{L_2}
=
\varepsilon.
\]

Similarly, for the Hessian channel,
\[
\sum_{t=0}^{T_j-1}\mathbb E\|H_t-\nabla^2F(x_t)\|_{\mathrm{op}}^3
\lesssim
\frac1{b_j^{3/2}}\sum_{t=0}^{T_j-1}\mathbb E\|s_t\|^3.
\]
Using \(\|s_t\|\lesssim \mu/L_2\) and \(\frac{T_j}{b_j^{3/2}}\sim n^{-1/4}2^{-5j/2}\),
\[
\|H_t-\nabla^2F(x_t)\|_{\mathrm{op}}
\lesssim
\mu.
\]
Choosing \(c_b,c_T\) appropriately and using that the homotopy depth \(J=\widetilde O(1)\) is logarithmic, the constants can be made smaller than any prescribed \(c_g,c_H\).
\end{proof}

\section{Proof of the local refinement lemma}\label{app:local}

\begin{proof}[Proof of \Cref{lem:local-main}]
We argue by induction over the stage iterates. Assume at time \(t\),
\[
\|\nabla F(x_t)\|\le A_g\varepsilon,
\qquad
\lambda_{\min}(\nabla^2F(x_t))\ge -A_H\mu.
\]
By \Cref{lem:pointwise-main},
\[
\|g_t-\nabla F(x_t)\|\le a_g\varepsilon,
\qquad
\|H_t-\nabla^2F(x_t)\|_{\mathrm{op}}\le a_H\mu.
\]
Hence
\[
\|g_t\|\le (A_g+a_g)\varepsilon,
\qquad
\lambda_{\min}(H_t)\ge -(A_H+a_H)\mu.
\]

Since \(m_t(s_t)\le 0\) and \(\Psi_{\beta_j}(s_t)\ge 0\),
\[
0
\ge
g_t^\top s_t+\frac12 s_t^\top H_t s_t+\frac M6\|s_t\|^3
\ge
-\|g_t\|\,\|s_t\|-\frac12 (A_H+a_H)\mu \|s_t\|^2+\frac M6\|s_t\|^3.
\]
Let \(r=\|s_t\|\). Then
\[
\frac M6 r^2-\frac12 (A_H+a_H)\mu r-(A_g+a_g)\varepsilon \le 0.
\]
Because \(\varepsilon=\mu^2/L_2\) and \(M=2L_2\), the positive root is \(O(\mu/L_2)\). Thus there exists \(C_R>0\) such that
\[
\|s_t\|\le C_R \frac{\mu}{L_2}.
\]

By choosing \(c_\beta\) sufficiently large, we ensure
\[
C_R\frac{\mu}{L_2}\le \rho_j=\frac{\beta_j}{M},
\]
so the entire stage lies in the small regime of \(\Psi_{\beta_j}\). Therefore \Cref{lem:psi-main} gives
\[
\|\nabla\Psi_{\beta_j}(s_t)\|\le M\|s_t\|^2,
\qquad
\|\nabla^2\Psi_{\beta_j}(s_t)\|_{\mathrm{op}}\le 2M\|s_t\|.
\]

Now use
\[
r_t=g_t+H_t s_t+\nabla\Psi_{\beta_j}(s_t)+\frac M2\|s_t\|s_t
\]
and
\[
\nabla F(x_{t+1})
=
\nabla F(x_t)+\nabla^2F(x_t)s_t+\zeta_t,
\qquad
\|\zeta_t\|\le \frac{L_2}{2}\|s_t\|^2.
\]
Exactly as in the proof of \Cref{lem:terminal-cert-main},
\[
\|\nabla F(x_{t+1})\|
\le
\|r_t\|+\|g_t-\nabla F(x_t)\|+\|H_t-\nabla^2F(x_t)\|_{\mathrm{op}}\|s_t\|+C L_2\|s_t\|^2.
\]
Every term on the right is \(O(\varepsilon)\), so after enlarging \(A_g\) if needed,
\[
\|\nabla F(x_{t+1})\|\le A_g\varepsilon.
\]

For the Hessian, the lower curvature certificate and \Cref{lem:psi-main} imply
\[
\lambda_{\min}(H_t)\ge -C M\|s_t\|.
\]
Hence,
\[
\lambda_{\min}(\nabla^2F(x_t))
\ge
\lambda_{\min}(H_t)-a_H\mu
\ge -C\mu.
\]
By \Cref{lem:hess-lip},
\[
\lambda_{\min}(\nabla^2F(x_{t+1}))
\ge
\lambda_{\min}(\nabla^2F(x_t))-L_2\|s_t\|
\ge -A_H\mu
\]
after enlarging \(A_H\) if necessary. This closes the induction.
\end{proof}

% \section{Implementation details}\label{sec:impl}

% \lstinputlisting[
%     style=pythonstyle,
%     caption={PyTorch implementation of the experimental pipeline.},
%     label={lst:exp-code}
% ]{re3mcn_experiment_fast.py}


% \newpage
% \input{checklist.tex}
\end{document}
